\section{\personBiran{}与\personCieliebak{}定理}
\label{sec:biran-cieliebak-theorem}

从定义看，一个几何上显然的事实在\personFloer{}同调不变量的应用中起关键作用；这就是
下面的消失结论。

\begin{theorem}
\label{thm:displaceable-floer-vanishes}
假设 $L\subset(M,\omega)$ 带有一个在哈密顿同痕下不变的“\personFloer{}同调”
$HF(L;M)$。若 $L$ 可位移，即存在哈密顿微分同胚 $\phi$，使
$\phi(L)\cap L=\varnothing$，则 $HF(L;M)=\{0\}$。
\end{theorem}

\begin{proof}
任取满足 $L\cap\phi(L)=\varnothing$ 的哈密顿微分同胚 $\phi$。按
$CF(L,\phi(L))$ 的定义，$CF(L,\phi(L))=\{0\}$，故
$HF(L,\phi(L))=\{0\}$。另一方面，哈密顿同痕不变性给出
$HF(L,\phi(L))\cong HF(L;M)$。证明完毕。
\end{proof}

回忆，我们已在第 16 章证明：若下列条件之一成立，则\personFloer{}同调良定义并在哈密顿
同痕下不变：
\begin{enumerate}[label=(\arabic*)]
\item $L_0,L_1$ 满足 $\Sigma_{L_i}>2$；
\item $L_1$ 与 $L_0=L$ 哈密顿同痕，且 $L$ 满足 $\Sigma_L\geq2$。
\end{enumerate}
特别地，对情形 (2)：
\begin{enumerate}[label=(\arabic*)]
\item 对任意满足 $L\pitchfork\phi(L)$ 的哈密顿微分同胚 $\phi$，
$HF(L,\phi(L))$ 良定义；
\item 若给定二者之间的哈密顿同痕，则 $HF(L,\phi(L))$ 与
$HF(L,\phi'(L))$ 之间存在典范同构，而且该同构只依赖于这一（哈密顿）同痕类。
\end{enumerate}

\begin{corollary}
\label{cor:monotone-lagrangian-R2n-HF-zero}
对 $\mathbb R^{2n}$ 中任意满足 $\Sigma_L\geq2$ 的紧单调拉格朗日子流形，
$HF(L,\phi(L))$ 具有上述性质，因而 $HF(L,\phi(L))=\{0\}$。
\end{corollary}

\personBiran{}与\personCieliebak{}在 (BCi02) 中观察到，紧拉格朗日嵌入的类似自动可位移性质在任意
如下类型的辛流形上都成立：
\begin{equation}
 V\times M,
\label{eq:17-subcritical-product}
\end{equation}
其中 $V$ 是次临界\personWeinstein{}流形，$M$ 是任意紧辛流形。特别地，他们以定理
17.3.9 和推论 17.4.2 结合这一观察，用一种有趣方式研究了紧拉格朗日子流形
$L\subset V\times\mathbb{CP}^n$ 的拓扑性质。本节余下部分紧随 (BCi02) 的论证
叙述其定理。

先从 (EG91) 引入非紧辛流形研究中的几个标准概念。

\begin{definition}
\label{def:convex-at-infinity}
若辛流形 $(M,\omega)$ 带有一个在无穷远处为完备辛膨胀的向量场 $X$，则称它\termConvexAtInfinity{}。称向量场 $X$ 完备辛膨胀，是指 $X$ 的流
$\{\phi^t\}$ 完备，且满足 $(\phi^t)^*\omega=e^t\omega$。我们假设
$(M,\omega)$ 在无穷远处容许一个多重次调和穷竭函数。
\end{definition}

设 $(Q,\xi)$ 是接触流形，$\lambda$ 是 $\xi$ 的接触形式，即满足
$\ker\lambda=\xi$ 的 1-形式。与 $\lambda$ 相伴的\personReeb{}向量场 $X_\lambda$ 是唯一
满足
\begin{equation}
 X_\lambda\mathbin\lrcorner\lambda=1,
 \qquad X_\lambda\mathbin\lrcorner d\lambda=0
\label{eq:17-reeb-definition}
\end{equation}
的向量场。因此接触形式自然给出分裂
\[
 TQ=\mathbb R\{X_\lambda\}\oplus\xi.
\]
以 $\Pi=\Pi_\lambda:TQ\to TQ$ 表示满足
$\operatorname{Image}\Pi=\xi$、
$\ker\Pi=\mathbb R\{X_\lambda\}$ 的幂等算子，并以
$\pi=\pi_\lambda:TQ\to\xi$ 表示相应投影。

下面是接触流形研究中的标准定义。

\begin{definition}
\label{def:strong-symplectic-filling}
若存在紧辛流形 $(W,\omega)$，满足
\begin{enumerate}[label=(\alph*)]
\item $\partial W=Q$；
\item 在边界附近，$\omega=d\lambda$，其中 $\lambda$ 是某个 1-形式，并且
$\partial W$ 上 $\xi=\ker\lambda$；
\item $\lambda\wedge(d\lambda)^{n-1}$ 所定义的 $Q$ 的定向，与
$Q=\partial W$ 的边界定向一致，
\end{enumerate}
\stfhIndex{I182}{Qiangxinketianchong}{强辛可填充（strongly symplectically fillable）}
则称紧接触流形 $(Q,\xi)$ 强辛可填充，并称 $(W,\omega)$ 为 $(Q,\xi)$ 的\termStrongSymplecticFilling{}。
\end{definition}

根据辛邻域定理，可在 $W$ 中 $\partial W=Q$ 的领形邻域 $U_\delta$ 上选择函数
$r$，使得
\begin{equation}
 U_\delta\cong(1-\delta,1]\times Q,
 \qquad \omega=d(r\pi^*\lambda)\quad\text{于 }U_\delta,
\label{eq:17-collar-contact-form}
\end{equation}
其中投影 $\pi:M^{\geq R}\to Q$。再考虑柱体
$(1-\delta,\infty)\times Q$，沿带状区域
$(1-\delta,1]\times Q\cong U_\delta$ 作并
\begin{equation}
 \widehat W=W\#((1-\delta,\infty)\times Q).
\label{eq:17-symplectic-completion}
\end{equation}
把 $\omega$ 与 $(1-\delta,\infty)\times Q$ 上的 $d(r\pi^*\lambda)$ 粘合，就自然
延拓到 $\widehat W=:M$。按定义，每个无穷远处凸的辛流形都有这种分解；只需令
$\varphi=\log r$（等价地 $r=e^\varphi$），并使用等同
\[
 M^{\geq R}=M^{r\geq R}\cong[0,\infty)\times\varphi^{-1}(\ln R).
\]
用 $Q$ 上选定的接触形式 $\lambda$ 和投影 $\pi:M^{\geq R}\to Q$，辛形式写为
\begin{equation}
 \widehat\omega=d(r\pi^*\lambda)=d(e^s\pi^*\lambda).
\label{eq:17-cylindrical-symplectic-form}
\end{equation}
称 $(s,y)$ 为\termCylindricalCoordinates{}，称 $(r=e^s,y)$ 为\termConeCoordinates{}。

\begin{remark}
\label{rem:cone-coordinates-examples}
当
$M=\mathbb C^n\setminus\{0\}\cong(0,\infty)\times S^{2n-1}$ 时，
$(\sqrt r,y)$（$y\in S^{2n-1}$）就是 $\mathbb C^n\setminus\{0\}$ 的标准极坐标。
另一方面，若写
\[
 T^*N\setminus\{0\}\cong S^1(T^*N)\times\mathbb R_+,
\]
则对 $T^*N\setminus\{0\}\subset T^*N$ 的典范坐标 $(q,p)$，有 $r=|p|$。
\end{remark}

在柱体 $(-\infty,0]\times Q\subset(-\infty,\infty)\times Q$ 上有自然分裂
\[
 T_xM\cong\mathbb R\left\{\frac1r\frac\partial{\partial r}\right\}
 \oplus T_qQ
 \cong\operatorname{span}\left\{\frac\partial{\partial s},X_\lambda(q)\right\}
 \oplus\xi_q\cong\mathbb R^2\oplus\xi_q.
\]
以 $\widetilde X_\lambda$ 表示 $M$ 上唯一的向量场，它在平移下不变，与 $r$ 的
水平集相切，并投影到 $X_\lambda$。不致混淆时，也把它简记为 $X_\lambda$。

下面描述一族适配于接触型 $(Q,\xi)$ 流形 $M$ 的给定柱面结构的特殊殆复结构。先在
配备接触形式 $\lambda$ 的接触流形 $(Q,\xi)$ 上定义 CR-殆复结构。

\begin{definition}[接触三元组与三元组度量]
\label{def:contact-triad-metric}
设自同态 $J_Q:TQ\to TQ$ 满足 $J_Q^2=-\Pi$。若双线性形式
$d\lambda(\cdot,J_Q\cdot)$ 在 $\xi$ 上非退化，则称 $J_Q$ 为 $\lambda$-相容
CR-殆复结构；称三元组 $(Q,\lambda,J_Q)$ 为\termContactTriad{}。称 $Q$ 上由
\[
 g_Q(h,k)=g_{(\lambda,J_Q)}:=\lambda(h)\lambda(k)+d\lambda(h,J_Qk),
 \qquad h,k\in TQ,
\]
定义的度量为\termTriadMetric{}。
\end{definition}

\begin{definition}
\label{def:lambda-contact-type-J}
设 $(M,\omega)$ 是如上的带柱面端辛流形。若在端上，殆复结构 $J$ 分裂为
\[
 J=j\oplus J_Q|_\xi:T M\cong\mathbb R^2\oplus\xi
 \longrightarrow TM\cong\mathbb R^2\oplus\xi,
\]
其中 $J_Q|_\xi$ 与 $d\lambda|_\xi$ 相容，且
$j:\mathbb R^2\to\mathbb R^2$ 把 $\partial/\partial s$ 映到 $X_\lambda$，则称
$J$ 为\termLambdaContactType{}。
\end{definition}

在 $M$ 上选择关于驯服殆复结构 $J$ 的多重次调和穷竭函数 $\varphi$。还假设在一个
紧集之外，$J$ 在 $X$ 的流下不变。于是，对充分大的 $S$，水平集
$\varphi^{-1}(S)$ 带有诱导接触结构，并且流映射
\[
 \phi:(0,\infty)\times\varphi^{-1}(S)
 \to M\cap\varphi^{-1}([S,\infty)),
 \qquad(s,y)\longmapsto\phi^s(y)
\]
定义一个微分同胚。在这些坐标中 $\varphi=s$，且 $\phi^s$ 就是平移映射
$\phi^s(S,y)=(S+s,y)$。由此，无穷远处凸辛流形的定义也可借助接触流形
$N=\varphi^{-1}(S)$ 给出。

\begin{definition}
\label{def:weinstein-subcritical}
若非紧辛流形 $(V,\omega)$ 上存在穷竭\personMorse{}函数
$\varphi:V\to\mathbb R$，以及梯度型的完备\personLiouville{}向量场 $X$（即满足
\stfhIndex{I198}{Wensitanliuxing}{温斯坦流形（Weinstein manifold）}
$\mathcal L_X\omega=\omega$），则称 $V$ 为\termWeinsteinManifold{}。若
$\varphi$ 在每个临界点处的\personMorse{}指标都严格小于
$n=\frac12\dim V$，则称 $V$ 为\termSubcritical{}；否则称为临界的。
\end{definition}

\begin{exercise}
\label{ex:unstable-manifold-isotropic}
设 $(V,\omega,\varphi)$ 如定义 17.4.8。证明任意
$p\in\operatorname{Crit}\varphi$ 的不稳定流形都是迷向的。
\end{exercise}

\begin{corollary}
\label{cor:psh-morse-index-bound}
多重次调和函数任意临界点的\personMorse{}指标不大于
$n=\frac12\dim V$。
\end{corollary}

从次临界\personWeinstein{}流形的如下普遍性质开始。

\begin{proposition}[{(BCi02) 引理 3.2}]
\label{prop:compact-set-displaceable-subcritical}
设 $(V,\omega,X,\varphi)$ 是完备次临界\personWeinstein{}流形，$A\subset V$ 是任意紧
子集。则存在紧支哈密顿量 $H$，使得 $\phi_H^1(A)\cap A=\varnothing$。
\end{proposition}

\begin{proof}
以 $\Delta_X$ 表示 $X$ 的流的（正向）不变集。按定义，它是不稳定流形的并；这个并
是紧的，并且是维数不大于 $n-1$ 的有限个光滑流形之并。因此可找出支于
$\Delta_X$ 的一个小邻域中的哈密顿量 $K$，使
$\phi_K^1(\Delta_X)\cap\Delta_X=\varnothing$。由于 $\Delta_X$ 紧，可找出小邻域
$\Delta_X\subset U$，使 $\phi_K^1(U)\cap U=\varnothing$。

以 $\psi^t$ 表示 $X$ 的流。因为 $\mathcal L_X\omega=\omega$，$\psi^t$ 是共形辛的，
故共轭 $\psi^t\circ\phi_K^1\circ(\psi^t)^{-1}$ 是辛的。记
$h_t=\psi^t\circ\phi_K^1\circ(\psi^t)^{-1}$。固定充分大的 $T>0$，使
$\psi^{-T}(A)\subset U$。因为 $\phi_K^1$ 把 $U$ 从自身移开，得到
$h_T(A)\cap A=\varnothing$。最后验证 $h_T$ 由下列 $H=H(t,x)$ 生成：
\[
 H(t,x)=e^{tT}(\psi^{-tT})^*(K_t+G_t),
\]
其中 $G_t$ 是由方程（见命题 3.4.8 的证明）
\[
 dG_t=\lambda_V-(\phi_K^t)^*\lambda_V
\]
唯一确定的紧支哈密顿量，并可显式写作
\[
 G(t,x)=\int_0^t(K_s-dK_s(X))((\phi_K^s)^{-1}(x))\,ds.
\]
由此，因 $K$ 紧支，$H$ 也紧支。
\end{proof}

\begin{lemma}
\label{lem:product-lagrangian-displaceable}
设 $(V,\omega_V)$ 是满足 $\omega_V=d\lambda_V$ 的次临界\personWeinstein{}流形，
$(M,\omega_M)$ 是闭辛流形。则任意紧拉格朗日子流形
$L\subset V\times M$ 都可由紧支哈密顿同痕位移。
\end{lemma}

\begin{proof}
用平行于 $TV\oplus\{0\}$ 的哈密顿向量场位移 $L$。由于 $L$ 紧，它在 $V$ 上的
投影也紧，结论立即由命题 17.4.11 得到。
\end{proof}

\stfhIndex{I041}{Bilan-Qilibake}{比兰--齐利巴克（Biran--Cielibak）}
\personBiran{}与\personCieliebak{} (BCi02) 随后证明了如下同调球面性结论。

\begin{theorem}[\personBiran{}--\personCieliebak{}]
\label{thm:biran-cieliebak-spherical}
设 $V$ 是满足 $2\leq\dim V\leq2n$ 且 $c_1(V)=0$ 的连通次临界\personWeinstein{}流形。
考虑任意满足 $H_1(L,\mathbb Z)=\{0\}$ 的紧拉格朗日子流形
$L\subset V\times\mathbb{CP}^n$。则 $L$ 是\termCohomologicallySpherical{}，
即对所有 $1\leq\ell\leq\dim L-1$，
\[
 H^\ell(L;\mathbb Z_2)=\{0\}.
\]
\end{theorem}

\begin{proof}
首先，由 $c_1(V)=0$ 且 $V$ 是\personWeinstein{}（特别地，它正合），
$V\times\mathbb{CP}^n$ 与 $\mathbb{CP}^n$ 有相同单调常数。我们知道
$\mathbb{CP}^n$ 是单调辛流形，即在
$\pi_2(M)\cong H_2(M;\mathbb Z)$ 上
$I_\omega=\lambda_0I_{c_1}$，其中
$\lambda_0=\omega(L)/c_1(L)>0$。
\footnote{译者注：本式中的 $L$ 在本章一直表示待研究的拉格朗日子流形，不能作为
$\omega$ 与 $c_1$ 的二维同调类输入；原书没有在此明确命名应取的投影直线类，故
译文照录而不作不可唯一确定的修补。}

接着证明，任意满足 $H_1(L,\mathbb Z)=\{0\}$ 的紧拉格朗日嵌入
$L\subset V\times\mathbb{CP}^n=:M$ 都单调，且
$\Sigma_L\geq\dim V=2n+1>2$。由于 $H_1(L,\mathbb Z)=\{0\}$，自然映射
\[
 H_2(V\times\mathbb{CP}^n;\mathbb Z)
 \longrightarrow H_2(V\times\mathbb{CP}^n,L;\mathbb Z)
\]
满射，故 $H_2(V\times\mathbb{CP}^n,L;\mathbb Z)$ 中每个元素都来自
$H_2(V\times\mathbb{CP}^n;\mathbb Z)$。
\footnote{译者注：本定理假设 $2\leq\dim V\leq2n$，但原书证明这里写
$\Sigma_L\geq\dim V=2n+1$，随后又写
$\Sigma_L=2(n+1)=\dim L+1$；这些维数关系对一般允许的 $\dim V$ 并不一致。
译文照录并登记，不替原书改定理条件。}

现在考虑 $\pi_2(M,L)$ 上的同态 $I_{(\omega;L)},I_{(\mu;L)}$。二者都经由
$H_2(M,L;\mathbb Z)$ 分解，并在其中的挠元上为零。另一方面，模挠元有
\[
 H_2(V\times\mathbb{CP}^n;\mathbb Z)
 \cong H_2(V;\mathbb Z)\oplus H_2(\mathbb{CP}^n;\mathbb Z).
\]
所以只需对 $H_2(V;\mathbb Z)$ 与 $H_2(\mathbb{CP}^n;\mathbb Z)$ 的生成元证明
\begin{equation}
 I_{(\omega;L)}=\frac{\lambda_0}{2}I_{(\mu;L)}.
\label{eq:17-relative-monotonicity}
\end{equation}
对 $H_2(V;\mathbb Z)$，由于 $\omega$ 正合且假设 $c_1(V)=0$，这显然成立。

令 $[C_0]\subset\mathbb{CP}^n$ 为复直线的同调类
$[C_0]\in H_2(\mathbb{CP}^n;\mathbb Z)$，并把它视为 $H_2(M;\mathbb Z)$ 的元素；
以 $i_*[C_0]$ 表示其到 $H_2(M,L;\mathbb Z)$ 的推前。显然
\[
 I_{(\omega;L)}(i_*[C_0])=\omega([C_0])=I_\omega([C_0]),
\]
\[
 I_{(\mu;L)}(i_*[C_0])=2c_1([C_0])=2I_{c_1}([C_0]).
\]
回忆 $2c_1([C_0])=2(n+1)>2$。此外，
$\Sigma_L=2(n+1)=\dim L+1$，故 $L$ 满足定理 17.3.9 (2) 的假设。这蕴含
\[
 H^\ell(L;\Lambda_{(\omega,L)})
 \cong HF^\ell(L,L;\Lambda_{(\omega,L)}),
 \qquad1\leq\ell\leq\dim L-1.
\]
另一方面，因为 $V$ 次临界，$V\times\mathbb{CP}^n$ 上存在紧支哈密顿微分同胚
$\phi$，使 $L\cap\phi(L)=\varnothing$，即 $L$ 可位移。结合定理 17.3.9 (2)
与推论 17.4.2，即证定理 17.4.13。
\end{proof}

下一节将结合定理 17.4.1 与第 17.3 节构造的谱序列，导出各种推论。
